\section{Introduction}
\label{sec:intro}

%% 目的：開場。三種合分方式擇一，加權平均最常見也最會誤導；立刻用 promise「公布 0.20、行為像 0.52」示範落差，
%%       並在此一次把 w_f、C_f 的記號交代掉（原本這兩個符號要到表格才第一次出現）。
A shared task with several dependent fields must decide how to combine them into one number, and
the common choice --- a weighted average of the per-field scores --- is the one whose weights
mislead. What an organizer publishes beside a field is not what that field is worth. In the NTCIR-19 AI CUP VeriPromise ESG track,
whose corpus is VeriPromiseESG4K~\cite{day2026veripromiseesg4k}, the promise field is published
at $0.20$ and, for our submission, behaved like $0.52$. Table~\ref{tab:schema} gives the four
fields, each field's published weight $w_f$ and scored class set $C_f$.

%% 目的：抽出造成落差的三個成分（固定權重的 macro-F1 平均、catch-all 類別、上游門控），據此推出兩個結論。
%%       注意舊版寫成「N/A 被排除在它所門控的欄位之外」，那句話文法上讓 N/A 變成門控者、事實上也和表格矛盾。
Three properties of that design open the gap, and none of them involves a model; the gap's width does. The per-field
scores are macro-F1, averaged at fixed weights. One class, here N/A, marks a field that carries
no real value; it is scored in some fields' averages and dropped from others. And an upstream
field \emph{gates} those below it: its value decides whether they carry a real class at all. Two
consequences follow once predictions respect the schema.

%% 目的：Lemma 1。macro-F1 對類別平均而非對樣本平均，所以 support 不進價格。用 val 上「1 筆 Misleading 對 566 筆
%%       Clear、同價」示範，再給一個讀者放得下尺度的比較。這裡的比較必須是完美 promise 門控（0.0770），
%%       不能是完美 binary（0.1565）——後者比 0.1167 大，寫成那樣會是假的。
\textbf{First}, macro-F1 averages over classes, not items. A scored class carries a share
$w_f/|C_f|$ of the composite score whatever its \emph{support}, the number of items carrying it;
how much of that share it collects depends on its own $F_1$, and on nothing else. Quality is
published at $0.35$ over three scored classes, so each of the three carries $0.1167$. One validation paragraph is labelled
\textsc{Misleading} --- greenwashing, where a company presents itself as greener than it is ---
and $566$ are labelled \textsc{Clear}; the two classes are worth exactly the same. That $0.1167$
is more than a perfect promise gate returns across all four fields ($0.0770$), and $3.7$ times
what a perfect \emph{binary} --- promise and evidence --- returns through the two fields it
gates ($0.0319$,
Section~\ref{sec:res:oracle}). Our system scores $F_1 = 0$ on \textsc{Misleading} on validation,
so the whole of that share is unclaimed.

%% 目的：Lemma 2。門控錯誤在兩個方向上定價不同、且方向相反；由此定義 effective weight（在第一次用到它的地方
%%       就定義，舊版是用完下一段才定義），給出 0.5216／0.4268，並補上「這是速率不是獎金」，否則讀者會拿它
%%       和後面的 0.0319 對撞。最後一次性地把「不是在批評評分設計」講完，避免同一句主張散在三處。
\textbf{Second}, a gate error is priced asymmetrically, and oppositely in its two directions. A
gate that wrongly shuts emits N/A, which in quality and timeline is not a scored class, so the
error costs recall and never precision; in evidence, where N/A is scored, the same error is
charged twice. A gate that wrongly opens is the mirror. A gate therefore carries not one
\emph{effective}
weight~\cite{wang1970differential} but two --- the derivative of the composite with respect to
that field's own macro-F1, one per direction, neither equal to its nominal $w_f$. Computing them
needs the published weights, the class distribution and one operating point --- the rates at
which the system in hand actually errs. Here the promise
field, published at $0.20$, is worth $0.5216$ when its gate recall improves and $0.4268$ when its
specificity does, because most of what it earns is collected in the fields it gates rather than
in itself. Neither consequence is a critique of the design: the weights are the organizers'
to set, and we leave them alone.

\begin{table}[t]
\caption{VeriPromiseESG4K. A class in parentheses is left out of that field's average;
N/A is scored in \emph{evidence} and excluded in \emph{quality} and \emph{timeline}. The
timeline classes are \texttt{already}, \texttt{within\_2\_years},
\texttt{between\_2\_and\_5\_years} and \texttt{more\_than\_5\_years} in the released files.}
\label{tab:schema}
\footnotesize
\setlength{\tabcolsep}{3pt}
\begin{tabular}{@{}lc>{\raggedright\arraybackslash}p{0.35\columnwidth}cc@{}}
\toprule
Field & $w_f$ & Labels & $|C_f|$ & $w_f/|C_f|$ \\
\midrule
promise (P)  & $0.20$ & Yes, No & 2 & $0.1000$ \\
evidence (E) & $0.30$ & Yes, No, N/A & 3 & $0.1000$ \\
quality (Q)  & $0.35$ & Clear, Not Clear, \textsc{Misleading}, (N/A) & 3 & $\mathbf{0.1167}$ \\
timeline (T) & $0.15$ & already, ${\le}2$y, 2--5y, ${>}5$y, (N/A) & 4 & $0.0375$ \\
\bottomrule
\end{tabular}
\end{table}

%% 目的：落到本資料集：語料規模、三個切分的大小（原本從未給，讓「三筆測試實例」無從放置），以及 schema 的兩條
%%       門控規則——用來說明級聯不是我們 pipeline 的產物。direct／cascade 的定義移到 §6.2 實際做分解的地方，
%%       因為在這裡定義完就要立刻破例，讀者得同時記住兩套互斥的分割。
VeriPromiseESG4K follows the ML-Promise schema~\cite{seki2025mlpromise} in Traditional Chinese:
four fields over $4{,}000$ paragraphs drawn from the sustainability reports of the fifty largest
listed Taiwanese companies, split $1{,}000$ train, $1{,}000$ validation and $2{,}000$ test. Two
schema rules do the gating: promise${=}$No forces evidence, quality and timeline to N/A, and
evidence$\in\{$No, N/A$\}$ forces quality to N/A. The gate is therefore the promise decision
itself, and by the second rule evidence is the gate on quality. Any schema-compliant prediction
inherits those two binary decisions~\cite{silla2011survey}: a team predicting timeline
``independently'' still answers the promise target with its N/A-versus-real split.

%% 目的：點出兩個結論指向相反的努力分配（追稀有類別 vs 修門控）——解決這個衝突才是本文要做的事。
The two consequences prescribe opposite allocations of effort: chase \textsc{Misleading}, or fix
the gate. \textsc{Misleading} is worth more, and is worth it whatever its support. The gate is
worth less, but unlike the rare class its gain yielded to ordinary engineering. Resolving
that conflict for our system, by measurement, is what this paper does.

%% 目的：定位相關工作的缺口（前人只講怎麼預測，沒講該往哪使力、也沒講支持數的份量），並交代分析所錨定的操作點。
Existing work reports how each field is predicted, not where the effort should go or what a
class's support is worth~\cite{wu2025cyut,nishi2025wcteam,leesombatwathana2025cscu}. Where effort
should go needs an operating point, and we use the one we submitted (Section~\ref{sec:system}).

%% 目的：C2：用數出來（而非擬合）的級聯模型，跨切分外推；把「值多少」與「修不修得到」分開；最後點出光是評分
%%       慣例本身造成的偏差就和任何建模增益一樣大。原本這段用了 roster、the scoring convention 等只出現一次
%%       又沒定義的名詞，改成讀者已經握有的詞。
\leadin{Contributions, by source of evidence.} \emph{Derived} from the scorer, the schema and
the label counts, before any prediction: a class's share of $S$ (Lemma~1), the gate written into
every downstream target (Lemma~2), and the form of a cascade model that turns error rates,
counted rather than fitted, into the four field scores. \emph{Measured on validation},
which is also our development set: the gate's effective weights, the price of each kind of
repair, a $-0.0265$ swing from how the composite is computed at all, larger than any gain our
ensemble or decision-layer work returned, and one deployed decision below what validation can
resolve (Section~\ref{sec:discussion}). \emph{On test, informed by leaderboard feedback} and so
not an independent estimate (Section~\ref{sec:setup}): the real-class gate-open rates counted on
validation hold there to within $0.05$, and of \textsc{Misleading}'s three instances we placed one, by search rather than
detection. Together these separate what an error is \emph{worth} from whether it can be
\emph{repaired}; we close with three additions to what a task overview reports.
